\subsection{MINORS OF A MATRIX}

Let \(X\) be a rectangular matrix \((\xi_{ij})_{(i,j)\in I\times J}\) of type \((I,J)\) whose indexing sets \(I\) and \(J\) are totally ordered. If \(H\subset I\) and \(K\subset J\) are finite subsets with the same number \(p\) of elements, there exists a unique increasing bijection \(\phi :H\to K\) (Set Theory, III, § 5, no. 3, Proposition 6); we shall denote by \(X_{H,K}\) the square matrix of type \((H,H)\) equal to \((\xi_{i,\phi (j)})_{(i,j)\in H\times H}\) . If the elements of \(X\) belong to a commutative ring \(A\) , the determinant \(\det (X_{H,K})\) is called the minor of the matrix \(X\) of indices \(H,K\) ; these determinants (for all ordered pairs \((H,K)\) of subsets of \(I\) and \(J\) respectively with \(p\) elements) are also called the minors of \(X\) of order \(p\) . With this notation:

PROPOSITION 9. Let M be an A-module with a basis \((e_i)_{i \in J}\) (finite or otherwise) whose indexing set J is totally ordered. For every integer \(p > 0\) , let \((e_H)_{H \in \mathfrak{F}_p(J)}\) be the corresponding basis of \(\bigwedge^p(M)\) (§ 7, no. 8). Let \((x_i)_{1 \leqslant i \leqslant p}\) be a sequence of \(p\) elements of M; let

\[
x _ {i} = \sum_ {j \in J} \xi_ {j i} e _ {j} \quad for i \in I = [1, p]
\]

and let \(X\) denote the matrix \((\xi_{ji})\) of type (J, I). Then

\[
x _ {1} \wedge x _ {2} \wedge \dots \wedge x _ {p} = \sum_ {\mathrm{H} \in \mathfrak {F} _ {p} (J)} (\det X _ {\mathrm{H,I}}) e _ {\mathrm{H}}.\tag{17}
\]

where H runs through the set \(\mathfrak{F}_{p}(\mathrm{J})\) of subsets of J with p elements.

This follows immediately from formula (12) of no. 4 and formula (6) of no. 3.

PROPOSITION 10. Let M and N be two free A-modules of respective dimensions m and n, \(u: \mathbf{M} \to \mathbf{N}\) a linear mapping and \(X\) the matrix of \(u\) with respect to a basis \((e_i)_{1 \leqslant i \leqslant m}\) of M and a basis \((f_j)_{1 \leqslant j \leqslant n}\) of N. Then, for every integer \(p \leqslant \inf(m, n)\) , the matrix of \(\bigwedge^p(u)\) with respect to the basis \((e_{\mathbf{K}})_{\mathbf{K} \in \mathfrak{F}_p(\mathbf{I})}\) of \(\bigwedge^p(\mathbf{M})\) and the basis \((f_{\mathrm{H}})_{\mathrm{H} \in \mathfrak{F}_p(\mathrm{J})}\) of \(\bigwedge^p(\mathbf{N})\) (where we have written \(\mathbf{I} = [1, m]\) and \(\mathbf{J} = [1, n]\)) is the matrix \((\det(X_{\mathrm{H},\mathbf{K}}))\) of type \((\mathfrak{F}_p(\mathbf{J}), \mathfrak{F}_p(\mathbf{I}))\) (and hence with \(\binom{n}{p}\) rows and \(\binom{m}{p}\) columns).

For a subset \(\mathbf{K} \subset \mathbf{J}\) with \(p\) elements, let \((j_k)_{1 \leqslant k \leqslant p}\) be the sequence of elements of \(\mathbf{K}\) arranged in increasing order; by definition of \(\bigwedge^p(u)\) , by §7, no. 2, formula (4)

\[
\bigwedge^ {p} (u) \left(e _ {\mathbf {K}}\right) = u \left(e _ {j _ {1}}\right) \wedge u \left(e _ {j _ {2}}\right) \wedge \dots \wedge u \left(e _ {j _ {p}}\right).
\]

Hence the element of the matrix of \(\bigwedge^{p}(u)\) which is in the row of index H and the column of index K is the component of index H of the element \(u(e_{j_1}) \wedge \cdots \wedge u(e_{j_p})\) ; it is therefore equal to \(\det(X_{\mathrm{H},\mathbf{K}})\) by Proposition 9.

The matrix \((\det(X_{\mathrm{H},\mathbf{K}}))\) is called the p-th exterior power of the matrix X and is denoted by \(\bigwedge^{p}(X)\) . When p = m = n, \(\bigwedge^{n}(X)\) is the matrix with the single element \(\det(X)\) .

PROPOSITION 11. Let M be a free A-module of finite dimension n; for every endomorphism u of M and every ordered pair of elements \(\xi\) , \(\eta\) of A,

\[
\det (\xi 1 _ {\mathrm{M}} + \eta u) = \sum_ {k \geqslant 0} \operatorname{Tr} \left(\bigwedge^ {k} (u)\right) \xi^ {n - k} \eta^ {k}.\tag{18}
\]

Let \((e_i)_{1\leqslant i\leqslant n}\) be a basis of M and let \(\mathbf{I} = [1,n]\) ; to calculate the left hand side of (18), we must form the product

\[
\left(\xi e _ {1} + \eta u (e _ {1})\right) \wedge \left(\xi e _ {2} + \eta u (e _ {2})\right) \wedge \dots \wedge \left(\xi e _ {n} + \eta u (e _ {n})\right)
\]

which is equal to the sum of the terms \(\xi^{n - p}\eta^{p}z_{\mathbf{K}}\) , where

\[
z _ {K} = x _ {1} \wedge x _ {2} \wedge \dots \wedge x _ {n}
\]

with \(x_{i} = u(e_{i})\) for \(i \in \mathbf{K}\) , \(x_{i} = e_{i}\) for \(i \in \mathbf{H} = \mathbf{I} - \mathbf{K}\) , where the integer \(p\) runs through the interval \([0, n]\) and, for each \(p, \mathbf{K}\) runs through the set of subsets of I with p elements. If \(i_{1} < i_{2} < \cdots < i_{n-p}\) (resp. \(j_{1} < j_{2} < \cdots < j_{p}\) ) are the elements of H (resp. K) arranged in increasing order, we can write (§7, no. 8, Corollary 1 to Theorem 1 and formula (19))

\[
z _ {\mathbf {K}} = \rho_ {\mathrm{H}, \mathbf {K}} e _ {i _ {1}} \wedge e _ {i _ {2}} \wedge \dots \wedge e _ {i _ {n - p}} \wedge u (e _ {j _ {1}}) \wedge \dots \wedge u (e _ {j _ {p}}).
\]

But if \(X\) is the matrix of \(u\) with respect to the basis \((e_i)\) , then by Proposition 10

\[
u (e _ {j _ {1}}) \wedge \dots \wedge u (e _ {j _ {p}}) = \sum_ {\mathrm{L} \in \mathfrak {F} _ {p} (\mathrm{I})} (\det X _ {\mathrm{L}, \mathrm{K}}) e _ {\mathrm{L}}
\]

and hence

\[
z _ {\mathbf {K}} = \rho_ {\mathrm{H}, \mathbf {K}} \sum_ {\mathbf {L} \in \mathfrak {F} _ {p} (\mathbf {I})} (\det X _ {\mathbf {L}, \mathbf {K}}) e _ {\mathrm{H}} \wedge e _ {\mathrm{L}}.
\]

Now \(H \cap L \neq \varnothing\) except for L = K; it therefore follows from §7, no. 8, formula (20) that \(z_{K} = (\det X_{K,K})e_{1} \wedge e_{2} \wedge \cdots \wedge e_{n}\) and formula (18) then follows from Proposition 10 and the definition of the trace of a matrix (II, §10, no. 11, formulae (49) and (50)).

COROLLARY. Under the same hypotheses as in Proposition 11, for the endomorphism \(\bigwedge(u)\) of the A-module \(\bigwedge(M)\)

\[
\operatorname{Tr} (\bigwedge (u)) = \det (1 _ {\mathbf {M}} + u).\tag{19}
\]

It suffices to replace \(\xi\) and \(\eta\) by 1 in (18) and observe that the matrix of \(\bigwedge(u)\) with respect to the basis of the \(e_{\mathbf{H}}\) ( \(\mathbf{H} \in \mathfrak{F}(\mathbf{I})\) ) is the diagonal matrix of the matrices of the \(\bigwedge^{k}(u)\) for \(k \geqslant 0\) (II, § 10, no. 7, Example IV).

